\appendix

\section*{Appendix X: Player Tools and Quick Reference}
\addcontentsline{toc}{section}{Appendix X: Player Tools and Quick Reference}
\label{app:player-tools}

This appendix provides lightweight, table‑friendly tools for tracking character progress, campaign state, and moment‑to‑moment rules. Print, photocopy, or copy to index cards for use at the table.

\subsection*{Physical Tracker Recommendations}
\label{app:physical-trackers}

Many groups find that physical, tactile trackers speed up play and reduce mental load. Consider:

\begin{itemize}
  \item \textbf{Poker chips or glass beads} for Boons and Story Beats (SB).  
  \emph{Keep a small bowl of chips in the center of the table – spend and earn visibly.}
  \item \textbf{Dry‑erase boards or laminated sheets} for clocks (4/6/8 segment circles).  
  \emph{Use a wet‑erase marker; draw a new clock each scene.}
  \item \textbf{Index cards (3×5)} for Assets and Followers.  
  \emph{Write the Asset’s name, free use, and surge on one side; flip when used.}
  \item \textbf{Colored paperclips or binder clips} on character sheets for Fatigue and Harm tracks.  
  \emph{Slide the clip up or down – no erasing.}
  \item \textbf{Printed “clock wheels”} (a circle divided into 4/6/8 wedges) placed inside a page protector.  
  \emph{Mark wedges with a dry‑erase marker, erase with tissue.}
  \item \textbf{D10 dice} to track temporary modifiers (Position, Effect, or Momentum).  
  \emph{Place the die next to the relevant tracker and tick up or down.}
\end{itemize}

\subsection*{Resolution Index Card}

\begin{tcolorbox}[colback=white,colframe=black,left=1em,right=1em,top=0.5em,bottom=0.5em]
\textbf{Declare Action → Build Dice Pool (Attribute + Skill) → Roll → Count Successes}

\medskip
0 Successes = \textbf{Miss} (GM gains SB, consequences) \\
1–3 Successes = \textbf{Partial} (success + complication) \\
4+ Successes = \textbf{Full Success} (clear win, no complication)

\smallskip
\textbf{Position} (Risk) + \textbf{Effect} (Impact) may shift outcomes.

\smallskip
\textbf{SB (Story Beats)} fuel abilities, magic, momentum, and narrative pushes.
\end{tcolorbox}

\subsection*{Fatigue and Harm Tracker}

\begin{tcolorbox}[colback=white,colframe=black,left=1em,right=1em,top=0.5em,bottom=0.5em]
\textbf{Fatigue:} \_\_\_\_\_\_\_\_  (Max = Body) \\
\textbf{Mental Fatigue:} \_\_\_\_\_\_\_\_ \\
\textbf{Harm:} [~] Harm 1 (–1 die) [~] Harm 2 (–2 dice) [~] Harm 3 (Incapacitated)

\medskip
When tracks fill:
\begin{itemize}
    \item Fatigue overflow → increase Harm by 1, clear Fatigue.
    \item Mental Fatigue → causes instability or disadvantage.
    \item Harm 3 → Last Stand or Patron’s Claim (see “Optional Player‑Character Death” in the Enhanced Player Play chapter).
\end{itemize}
\end{tcolorbox}

\subsection*{SB and Momentum Tracking}

\begin{tcolorbox}[colback=white,colframe=black,left=1em,right=1em,top=0.5em,bottom=0.5em]
\textbf{Current SB (GM pool):} \_\_\_\_\_\_ \\
\textbf{Momentum (party pool):} \_\_\_\_\_\_  (max 2) \\
\textbf{Used this session:} \_\_\_\_\_\_
\end{tcolorbox}

\subsection*{Talent Progression}

\begin{center}
\begin{tabular}{|c|c|}
\hline
\textbf{XP Cost} & \textbf{Talent Type / Impact} \\
\hline
2 XP & Minor Talent / Perk upgrade / situational bonus \\
4 XP & Major Talent / specialization / strong permanent effect \\
6 XP & Access talent (Spellcraft, Familiar, Codex, Patron’s Symbol) \\
8 XP & Prestige Talent / campaign‑defining ability \\
\hline
\end{tabular}
\end{center}

\subsection*{Character Goals Card}

\begin{tcolorbox}[colback=white,colframe=black,width=\textwidth,left=1em,right=1em]
\textbf{Name:} \_\_\_\_\_\_\_\_\_\_\_\_ \\
\textbf{Core Drive:} \_\_\_\_\_\_\_\_\_\_\_\_ \\
\textbf{Short-Term Goal (1–2 sessions):} \_\_\_\_\_\_\_\_\_\_\_\_ \\
\textbf{Long-Term Goal (arc):} \_\_\_\_\_\_\_\_\_\_\_\_ \\
\textbf{Background Complication:} \_\_\_\_\_\_\_\_\_\_\_\_ \\
\textbf{Relationship That Matters:} \_\_\_\_\_\_\_\_\_\_\_\_
\end{tcolorbox}

\emph{Physical tip:} Print this card on cardstock, fold in half, and stand it in front of you as a reminder.

\subsection*{Faction Attitude Card}

\begin{tcolorbox}[colback=white,colframe=black,left=1em,right=1em,top=0.5em,bottom=0.5em]
\textbf{Faction Name:} \_\_\_\_\_\_\_\_\_ \\
Attitude: –3 –2 –1 0 +1 +2 +3 \\
Notes: \_\_\_\_\_\_\_\_\_\_\_\_\_\_\_\_\_\_\_
\end{tcolorbox}

\subsection*{Player Orientation Card (Quick Principles)}

\begin{tcolorbox}[colback=white,colframe=black,left=1em,right=1em]
\textbf{1. Declare Bold Action} — Drive fiction forward. \\
\textbf{2. Accept Consequences} — Complications create story. \\
\textbf{3. Spend Story Beats} — Make moments matter. \\
\textbf{4. Use Boons} — Re‑roll, improve Position, activate Assets. \\
\textbf{5. Share the Spotlight} — Assist, bond, set up allies.
\end{tcolorbox}

\subsection*{Class / Archetype Resource Quick Reference}

\begin{center}
\begin{tabular}{|l|l|}
\hline
\textbf{Archetype} & \textbf{Core Resource} \\
\hline
Fighter / Warrior & Maneuvers, Fatigue, Weapon Mastery \\
Rogue / Scout & Positioning, Stealth, Backstab \\
Caster (Freeform) & Spellcraft, Elemental Tags, Backlash \\
Runekeeper & Rites, Obligation, Codex \\
Invoker & Rituals, Patron’s Symbol, Crack the Seal \\
Cantor & Songs, Corruption Clock, Repertoire \\
Summoner & Leash Clock, Spirit Bond, Pact \\
\hline
\end{tabular}
\end{center}

\subsection*{How to Engage with the GM (Player Prompts)}

\begin{itemize}
    \item \textbf{Ask:} “What do I need to roll?” – get the DV and Position.
    \item \textbf{Propose:} “Can I spend SB to achieve X?” – negotiate narrative pushes.
    \item \textbf{Declare:} “Here’s my intent and method.” – always state what you want to happen.
    \item \textbf{Offer:} “I’ll take a complication if Y happens.” – invite Devil’s Bargains.
    \item \textbf{Clarify:} “What’s at risk if I fail?” – know the stakes before rolling.
\end{itemize}

\subsection*{Downtime Project Tracker}

\begin{tcolorbox}[colback=white,colframe=black,left=1em,right=1em]
\textbf{Name:} \_\_\_\_\_\_\_\_\_\_\_  \qquad \textbf{Segments:} [ ] [ ] [ ] [ ] [ ] [ ] \\
\textbf{Goal:} \_\_\_\_\_\_\_\_\_\_\_ \\
\textbf{Obstacles / Resources Needed:} \_\_\_\_\_\_\_\_
\end{tcolorbox}

\emph{Physical tip:} Use a 6‑segment whiteboard clock. Each downtime, erase one segment and roll for progress.

\subsection*{What Consequence Do You Accept? (Menu for GM)}
\label{app:consequence-menu}

When a player fails or partially succeeds, the GM may offer a choice from this list (or invent a custom one):
\begin{itemize}
    \item Harm (level 1 or 2)
    \item Fatigue (1 or 2)
    \item Lost Advantage (Position worsens, a clock advances)
    \item Time Delays (a clock ticks, a deadline tightens)
    \item Exposure / Heat (attention from enemies, bounty)
    \item Resource Depletion (supply clock tick, gear compromised)
    \item Relationship Strain (a bond or faction attitude worsens)
\end{itemize}

\subsection*{Ways to Shine in Fate's Edge}

\begin{itemize}
    \item Tie actions to your drives and backstory – invoke your Background Boon.
    \item Offer complications and accept them – earn Boons.
    \item Spend SB to escalate or transform the scene – make failure memorable.
    \item Use your class resource early and often – don’t hoard.
    \item Connect scenes to campaign themes – name the faction or place.
    \item Make other characters look cool – assist, set up, pass the spotlight.
\end{itemize}

\subsection*{Blank Player Index Cards (for quick notes)}

\begin{center}
\fbox{\parbox{0.75\textwidth}{\vspace{0.75in} \centering Name / Boons / Fatigue / Harm}}
\end{center}

\begin{center}
\fbox{\parbox{0.75\textwidth}{\vspace{0.75in} \centering Clocks / Obligation / Scene Notes}}
\end{center}

\subsection*{Final Reminder}

\textbf{Your character sheet is not your character.}  
It is a tool to help you express who they have been, who they are becoming, and how their choices change the world.

Use these tools to make that expression vivid, consequential, and memorable.

\clearpage

\section*{Appendix Y: Magic Systems Reference}
\addcontentsline{toc}{section}{Appendix Y: Magic Systems Reference}
\label{app:magic-ref}

Magic in Fate’s Edge is plural.  
Different traditions channel power through different metaphysical structures, but all obey the same truth:

\begin{center}
\emph{Power has a cost, and magic shapes the one who wields it.}
\end{center}

This appendix provides quick‑reference tools for players using magic‑affiliated systems, talismans, rites, and occult practices.

\subsection*{Core Magical Principles}
\begin{tcolorbox}[colback=white,colframe=black,left=1em,right=1em]
\begin{itemize}[noitemsep]
    \item Magic expresses \textbf{identity and belief}.
    \item Casting always creates \textbf{consequences somewhere}.
    \item Magic rewrites reality through \textbf{intention, structure, or sacrifice}.
    \item Excessive power erodes \textbf{choice} (see Iconic Ascension sidebar in the full rules).
\end{itemize}
\end{tcolorbox}

\subsection*{Universal Casting Reference}
\begin{tcolorbox}[colback=gray!5,colframe=gray!60]
\textbf{Cast Declaration Template:}
\begin{enumerate}[noitemsep]
    \item \textbf{Intent} — What change do you seek?
    \item \textbf{Shape} — How does your tradition express it?
    \item \textbf{Cost} — What do you risk or expend?
    \item \textbf{Resolution} — Test; apply Position/Effect rules.
    \item \textbf{Backwash} — Mark Fatigue, SB pressure, taboo, instability, or bargain.
\end{enumerate}
\end{tcolorbox}

\subsection*{Magic Tradition Quick Reference}
\label{app:magic-traditions}

\begin{description}
  \item[Freeform Arcana (Wizardry)] \textit{Creative, rational, deliberate shaping.} Tools: Glyphs, structures, metaphors. Costs: Fatigue, reality recoil, SB instability. \textbf{Play Cue:} Diagram your spell or state its symbolism. Big spells distort fate – accept SB spikes.
  
  \item[Innate Arcana (Sorcery)] \textit{Emotional, instinctive, bloodline‑linked.} Tools: Breath, impulse, resonance. Costs: Fatigue, unpredictable backlash. \textbf{Play Cue:} Treat magic like a reflex. Push your limits – expect consequences.
  
  \item[Invocation (Priest / Paladin / Medium)] \textit{Petition, prayer, pact, rite.} Tools: Symbols, vows, places of power. Costs: Taboo pressures, obligation clocks. \textbf{Play Cue:} Your magic is not your own – frame it as a request. Violating your ethos draws consequences.
  
  \item[Runic Consecration (Runekeeping)] \textit{Methods, seals, contracts.} Tools: Ritual steps, codices, binding clauses. Costs: Hunger clocks, escalation loops. \textbf{Play Cue:} Magic is \emph{procedure}. Expect the system to push back.
  
  \item[Shaping / Druidic Power] \textit{Resonance with place, element, cycle.} Tools: Marks, offerings, seasonal rites. Costs: Environmental echo, territorial consequence. \textbf{Play Cue:} Declare what the land wants – then bargain.
  
  \item[Pact Binding (Warlock / Avatar)] \textit{Exchange, hunger, transformation.} Tools: Tokens, signatures, reflection marks. Costs: Autonomy loss, ascension risk. \textbf{Play Cue:} Ask: “What does my patron get?” Debt becomes plot.
  
  \item[Summoning (Medium / Necromancer / Shaman)] \textit{Invocation of entities, binding, vassal action.} Tools: Circles, bargains, lanterns, ancestral calls. Costs: Attention clocks, breach risk. \textbf{Play Cue:} Give the summoned a personality – their aid comes with demands.
\end{description}

\subsection*{Mutual Rules of Magical Cost}
\begin{center}
\textbf{No magic is free.}
\end{center}

\begin{center}
\begin{tabular}{|c|p{9cm}|}
\hline
\textbf{Cost Type} & \textbf{Examples in Play} \\
\hline
Fatigue & Exhaustion surge, shaking focus, body strain \\
SB Pressure & GM spikes complications or clocks (elemental backlash, Patron omens) \\
Taboo / Obligation & Commandments, rites, punishment for misuse, Patron demands \\
Attention / Backwash & Spirits notice; cosmic ledger ticks; faction awareness \\
Anomaly / Distortion & Places react, reality warps, rituals echo – scene tags change \\
Identity Erosion & Compelled acts, alignment tilt, Iconic Ascension risk \\
\hline
\end{tabular}
\end{center}

\subsection*{Cross‑System Synergies}
All magic systems can blend. A Runekeeper may emulate a Cleric or Druid; an Invoker may operate as a Paladin or Prophet; a Sorcerer may look like a Witch, Elementalist, or Psion.

When combining systems:
\begin{enumerate}
    \item Choose your \textbf{primary expression} (belief, intellect, instinct, ritual).
    \item Describe spells through that lens.
    \item Let costs reflect your system.
    \item Track consequences separately – mixing systems increases risk.
\end{enumerate}

\subsection*{GM Reminders for Magic}
\begin{itemize}
    \item Magic creates story – reward its use with consequences and momentum.
    \item Players should narrate method, not just effect.
    \item Big moments should change scenes, not just numbers.
    \item Magic should tempt the caster – echo their drives and flaws.
\end{itemize}

\subsection*{Player Reminders for Magic}
\begin{itemize}
    \item Lean into cost – magic feels significant when it hurts.
    \item Make magic personal – describe what it looks and feels like.
    \item Tie powers to identity, belief, trauma, duty, or legacy.
    \item Let magic change your character over time.
\end{itemize}

\clearpage

\section*{Appendix Z: Pre‑made Spells with TAGS and DV Math}
\addcontentsline{toc}{section}{Appendix Z: Pre‑made Spells with TAGS and DV Math}
\label{app:spell-examples}

The following examples illustrate how TAGS, DV scaling, and spell intent interact within the casting system. Use them as inspiration or as templates for your own spells.

\subsection*{How to Read a Spell Entry}
Each entry includes:
\begin{itemize}
  \item \textbf{Name} and \textbf{Tradition} (optional).
  \item \textbf{TAGS} – the components that define the effect.
  \item \textbf{Effect} – plain‑language outcome.
  \item \textbf{DV} – base calculation (1 + number of TAGS; dangerous tags add +2).
  \item \textbf{Notes} – typical range, duration, or special considerations.
\end{itemize}

\subsection*{Tier I: Minor Utility (DV 2–3)}

\paragraph{Simple Light} (Air + Light) \\
\textbf{TAGS:} {[ILLUMINATE]} \\
\textbf{Effect:} Create a small sustained light (torch‑brightness) for the scene. \\
\textbf{DV:} 1 + 1 = \textbf{2}. \\
\textbf{Notes:} No backlash risk; lasts until dismissed.

\paragraph{Mend Minor Wounds} (Life + Mend) \\
\textbf{TAGS:} {[MEND]} \\
\textbf{Effect:} Close a small cut or bruise; reduce 1 Fatigue (Minor). \\
\textbf{DV:} 1 + 1 = \textbf{2}. \\
\textbf{Notes:} Requires touch.

\begin{tcolorbox}[
    colback=parchment,
    colframe=bronze,
    sharp corners,
    breakable,
    enhanced,
    title=Healing Harm and Fatigue,
    fonttitle=\bfseries,
    attach boxed title to top left={yshift=-2mm,xshift=4mm},
    boxed title style={colback=bronze, colframe=bronze, sharp corners, rounded corners=north}
]

\textbf{Mend as a Cantrip.} The \texttt{Mend} tag (Life) can be cast as a single‑action cantrip (DV 2). On a success, you remove \textbf{1 Fatigue} from yourself or a touched ally. This cannot heal Harm.

\textbf{Healing Harm (1 level).} To restore a level of Harm, you must include the \texttt{Mend} tag and either:
\begin{itemize}
    \item add a \textbf{+2 DV modifier} to the spell, or
    \item combine \texttt{Mend} with a second appropriate tag (e.g., \texttt{Purify}, \texttt{Strengthen}, \texttt{Cleanse}).
\end{itemize}
The GM sets the final DV based on the tags and the wound’s severity.

\textit{Why?} Fatigue is exhaustion – a quick charm can ease tired muscles. Harm is injury – closing wounds demands more power and risk. This keeps healing meaningful without overshadowing martial endurance.

\end{tcolorbox}

\paragraph{Message on Wind} (Air + Veil) \\
\textbf{TAGS:} {[WHISPER]} \\
\textbf{Effect:} Carry a short whispered message 100 feet on wind currents. \\
\textbf{DV:} 1 + 1 = \textbf{2}. \\
\textbf{Notes:} One‑way communication; 30 seconds duration.

\subsection*{Tier II: Combat Utility (DV 3–4)}

\paragraph{Blade Flame} (Fire + Strike) \\
\textbf{TAGS:} {[BURNING][STRIKE]} \\
\textbf{Effect:} Coat a weapon in fire, adding +1 Harm and {[BURN]} to the next melee attack. \\
\textbf{DV:} 1 + 1 + 1 = \textbf{3}. \\
\textbf{Notes:} Lasts one exchange.

\paragraph{Shield of Air} (Air + Wall) \\
\textbf{TAGS:} {[WARD][DEFEND]} \\
\textbf{Effect:} Create a compressed air barrier; gain Position +1 on next defense. \\
\textbf{DV:} 1 + 1 + 1 = \textbf{3}. \\
\textbf{Notes:} One use only.

\paragraph{Fleet Step} (Air + Leap) \\
\textbf{TAGS:} {[HASTE][MOVE]} \\
\textbf{Effect:} Double movement speed for one exchange. \\
\textbf{DV:} 1 + 1 + 1 = \textbf{3}. \\
\textbf{Notes:} Requires open area.

\subsection*{Tier III: Area Control (DV 4–5)}

\paragraph{Wall of Stone} (Earth + Wall) \\
\textbf{TAGS:} {[BARRIER][FORTIFY]} \\
\textbf{Effect:} Summon a 10‑foot stone wall (Near; scene duration). \\
\textbf{DV:} 1 + 1 + 1 = \textbf{3}. \emph{Dangerous DV:} \textbf{4} (combat blocking). \\
\textbf{Notes:} Provides cover; obstructs movement.

\paragraph{Sleep’s Gentle Touch} (Water + Mind) \\
\textbf{TAGS:} {[DREAM][COMMAND]} \\
\textbf{Effect:} One target falls asleep (scene duration; Resist: Spirit + Resolve DV 2). \\
\textbf{DV:} 1 + 1 + 1 = \textbf{3}. \emph{Dangerous DV:} \textbf{4} (combat use). \\
\textbf{Notes:} Harm wakes target.

\paragraph{Flame Cloak} (Fire + Area) \\
\textbf{TAGS:} {[BURNING][AREA]} \\
\textbf{Effect:} 6‑foot radius of intense heat; enemies suffer –1 die to approach. \\
\textbf{DV:} 1 + 1 + 1 = \textbf{3}. \emph{Dangerous DV:} \textbf{4} (if allies are near). \\
\textbf{Notes:} Creates difficult terrain.

\subsection*{Tier IV: Battlefield Shaping (DV 5–6)}

\paragraph{Earth’s Tremor} (Earth + Force) \\
\textbf{TAGS:} {[QUAKE][AREA]} \\
\textbf{Effect:} 20‑foot tremor; enemies test Body + Athletics (DV 3) or fall prone. \\
\textbf{DV:} 1 + 1 + 1 = \textbf{3}. \emph{Dangerous DV:} \textbf{5}. \\
\textbf{Notes:} Destabilizes structures; affects all in area.

\paragraph{Storm’s Fury} (Air + Storm) \\
\textbf{TAGS:} {[STORM][AREA][FORCE]} \\
\textbf{Effect:} 30‑foot storm cloud; lightning strikes random targets (Harm 2). \\
\textbf{DV:} 1 + 1 + 1 + 1 = \textbf{4}. \emph{Dangerous DV:} \textbf{6}. \\
\textbf{Notes:} Lasts scene; reduces visibility.

\paragraph{Healing Circle} (Life + Area) \\
\textbf{TAGS:} {[MEND][AREA][CLEANSE]} \\
\textbf{Effect:} Allies in 15‑foot radius heal 1 Harm and clear 1 Fatigue. \\
\textbf{DV:} 1 + 1 + 1 + 1 = \textbf{4}. \\
\textbf{Notes:} Major positive effect warrants high DV.

\subsection*{Tier V: Domain‑Level Effects (DV 6–7)}

\paragraph{Gate of Passage} (Air + Fold) \\
\textbf{TAGS:} {[GATE][TRANSPORT][AREA]} \\
\textbf{Effect:} Open a two‑way portal 20 feet wide linking two locations. \\
\textbf{DV:} 1 + 2 + 1 + 1 = \textbf{5}. \emph{Dangerous DV:} \textbf{7}. \\
\textbf{Notes:} Requires concentration; may destabilize.

\paragraph{Army of Stone} (Earth + Summon) \\
\textbf{TAGS:} {[CONJURE][AREA][GUARD]} \\
\textbf{Effect:} Summon 4 stone guardians (Cap 2) to follow simple commands. \\
\textbf{DV:} 1 + 2 + 1 + 1 = \textbf{5}. \emph{Dangerous DV:} \textbf{6}. \\
\textbf{Notes:} Scene duration; concentration.

\paragraph{Mind’s Dominion} (Fate + Mind) \\
\textbf{TAGS:} {[COMMAND][AREA][CURSE]} \\
\textbf{Effect:} All enemies in 25‑foot radius test Spirit + Resolve (DV 4) or obey a single command. \\
\textbf{DV:} 1 + 1 + 1 + 2 = \textbf{5}. \emph{Dangerous DV:} \textbf{7}. \\
\textbf{Notes:} Significant ethical risk; backlash likely.

\subsection*{Signature / Personal Spells (Custom)}

\paragraph{The Warden’s Bond} (Life + Bind) – Healing Cantor \\
\textbf{TAGS:} {[MEND][BIND][GUARD]} \\
\textbf{Effect:} Heal 2 Harm and link caster to target – the next Harm they take, caster suffers 1 instead. \\
\textbf{DV:} 1 + 1 + 1 + 1 = \textbf{4}. \\
\textbf{Notes:} Sacrificial protection.

\paragraph{Shadow’s Step} (Air + Veil) – Trickster’s Path \\
\textbf{TAGS:} {[VEIL][MOVE][LEAP]} \\
\textbf{Effect:} Become invisible and teleport 30 feet; gain +2 dice to next Stealth action. \\
\textbf{DV:} 1 + 1 + 1 + 1 = \textbf{4}. \\
\textbf{Notes:} Requires shadow or dim light.

\paragraph{Forge’s Truth} (Fire + Reveal) – Artificer’s Craft \\
\textbf{TAGS:} {[REVEAL][MEND][FORTIFY]} \\
\textbf{Effect:} See flaws in an object and repair as if 2 Skill levels higher (scene). \\
\textbf{DV:} 1 + 1 + 1 + 1 = \textbf{4}. \\
\textbf{Notes:} Works on magical items (with risk).

\subsection*{Backlash and Push Guidelines}

\subsubsection*{Backlash Examples by Tier}
\begin{itemize}
    \item \textbf{Tier I–II:} Minor fatigue, brief elemental effect (smoke, chill, static).
    \item \textbf{Tier III–IV:} Harm 1, temporary Condition, environmental hazard.
    \item \textbf{Tier V:} Harm 2, major Condition, significant environmental change.
\end{itemize}

\subsubsection*{Push Options (Accelerate or Amplify)}
\begin{itemize}
    \item \textbf{+1 Effect:} Increase radius, duration, or potency one step.
    \item \textbf{+1 die:} Improve casting roll (cost: +1 Fatigue).
    \item \textbf{Instant:} Resolve immediately (cost: +1 Harm).
    \item \textbf{Sustain:} Maintain effect with concentration (+1 Fatigue/round).
\end{itemize}

\noindent This framework provides concrete spell examples and demonstrates how TAGS, DV scaling, and ambition interact within the magic system.